\subsection{Functions of positive type on a topological group}

In this number, G is a topological group whose identity element is denoted by e. The Hilbert spaces considered are complex.

DEFINITION 4. --- A continuous function \(\varphi\in\mathcal{C}(\mathrm{G})\) is said to be of positive type on G if the function \(f\) defined by \(f(g,h)=\varphi(g^{-1}h)\) on \(\mathrm{G}\times\mathrm{G}\) is a universally positive kernel on \(\mathrm{G}\) .

We denote by \(\operatorname{Pos}(\mathbf{G})\) the set of functions of positive type on G and we denote by \(\operatorname{Pos}_{1}(\mathbf{G})\) (resp. \(\operatorname{Pos}_{\leqslant1}(\mathbf{G})\) ) the subset of \(\varphi\in\operatorname{Pos}(\mathbf{G})\) such that \(\varphi(e)=1\) (resp. such that \(\varphi(e)\leqslant1\) ).

Example. --- Let \(\varrho\) be a unitary representation of G in a Hilbert space E and \(x \in \mathrm{E}\) . Let \(\varphi\) be the diagonal matrix coefficient defined by \(\varphi(g) = \langle x \mid \varrho(g)x \rangle\) for all \(g \in \mathrm{G}\) ; the function \(\varphi\) is continuous. It follows that

\[
\varphi (g ^ {- 1} h) = \langle x | \varrho (g) ^ {*} \varrho (h) x \rangle = \langle \varrho (g) x | \varrho (h) x \rangle
\]

for all \((g,h)\in\mathrm{G}\times\mathrm{G}\) . Consequently \(\varphi\in\mathrm{Pos}(\mathrm{G})\) (theorem 1 of V, p. 432, (iv)). We have \(\varphi\in\mathrm{Pos}_{1}(\mathrm{G})\) if and only if \(\|x\|=1\) .

DEFINITION 5. --- Let \(\varphi\) be a function of positive type on G. A Hilbert realization of \(\varphi\) is a pair \((\varrho, x)\) where \(\varrho\) is a unitary representation of G in a Hilbert space E and \(x \in \mathrm{E}\) such that \(\varphi(g) = \langle x \mid \varrho(g)x \rangle\) for all \(g \in \mathrm{G}\) .

If x is a cyclic vector of \(\varrho\) , one says that it is a cyclic Hilbert realization of \(\varphi\) .

PROPOSITION 9. --- Let \(\varphi\in\mathcal{C}(\mathrm{G})\) . The following conditions are equivalent:

\begin{enumerate}
\def\labelenumi{(\roman{enumi})}
\item
  The function \(\varphi\) is of positive type on G;
\item
  There exists a cyclic Hilbert realization of \(\varphi\) ;
\item
  There exists a Hilbert realization of \(\varphi\) .
\end{enumerate}

Condition (ii) implies condition (iii), and condition (iii) implies (i) by the example above.

Let us finally prove that (i) implies (ii). Let (E, \(\gamma\) ) be a cyclic Hilbert realization of the universally positive kernel defined by \(f(g,h)=\varphi(g^{-1}h)\) for \((g,h)\in\mathrm{G}\times\mathrm{G}\) . Let \(k\in\mathrm{G}\) . The continuous function \(\gamma_{k}:g\mapsto\gamma(kg)\) on G satisfies

\[
\langle \gamma_ {k} (g) | \gamma_ {k} (h) \rangle = \langle \gamma (k g) | \gamma (k h) \rangle = f (k g, k h) = \varphi ((k h) ^ {- 1} k g) = f (g, h)
\]

for all \((g,h)\in\mathrm{G}\times\mathrm{G}\) , hence \((\mathrm{E},\gamma_{k})\) is a Hilbert realization of f. By prop. 1 of V, p. 434, there exists a unique unitary operator \(\varrho(k)\) in \(\mathcal{L}(\mathrm{E})\) such that \(\gamma_{k}=\varrho(k)\circ\gamma\) . For every \(g\in G\) and every \((k_{1},k_{2})\in\mathrm{G}\times\mathrm{G}\) , it follows that

\[
\varrho (k _ {1} k _ {2}) (\gamma (g)) = \gamma (k _ {1} k _ {2} g) = \varrho (k _ {1}) (\varrho (k _ {2}) (\gamma (g))),
\]

whence \(\varrho(k_1k_2) = \varrho(k_1)\varrho(k_2)\) since the image of \(\gamma\) is a total subset of E.

Let \(k \in G\) . For every \(g \in G\) , we have \(\varrho(g)(\gamma(k)) = \gamma(gk)\) , and the map from G to E defined by \(g \mapsto \varrho(g)(\gamma(k))\) is therefore continuous. Since the endomorphism \(\varrho(g)\) is unitary, one deduces that \(\varrho\) is a unitary representation of G in E (lemma 4 of V, p. 380). Then set \(x = \gamma(e)\) . The set of vectors \(\varrho(g)x = \varrho(g)(\gamma(e)) = \gamma(g)\) for \(g\) ranging over G is total in E, hence x is a cyclic vector of \(\varrho\) . Since

\[
\langle x \mid \varrho (g) x \rangle = f (e, g) = \varphi (g)
\]

for all \(g \in G\) , the pair \((\varrho, x)\) is a cyclic Hilbert realization of \(\varphi\) .

PROPOSITION 10. — Let \(\varphi\) be a function of positive type on G and let \((\varrho_{1}, x_{1})\) and \((\varrho_{2}, x_{2})\) be Hilbert realizations of \(\varphi\) , the Hilbert representation \((\varrho_{1}, x_{1})\) being cyclic. There exists a unique isometric G-morphism \(u\) from \(\varrho_{1}\) into \(\varrho_{2}\) such that \(u(x_{1}) = x_{2}\) . If \((\varrho_{2}, x_{2})\) is also cyclic, then \(u\) is an isomorphism.

For \(1 \leqslant j \leqslant 2\) , let \(\mathrm{E}_j\) be the space of \(\varrho_j\) and \(\gamma_j\) the function on \(\mathbf{G}\) defined by \(\gamma_j(g) = \varrho_j(g)x_j\) for all \(g \in \mathbf{G}\) . The pairs \((\mathrm{E}_1, \gamma_1)\) and \((\mathrm{E}_2, \gamma_2)\) are Hilbertian realizations of the function of positive type \((g, h) \mapsto \varphi(g^{-1}h)\) , and \((\mathrm{E}_1, \gamma_1)\) is a cyclic Hilbertian realization. By prop. 1 of V, p. 434, there exists a unique isometric linear map \(u: \mathrm{E}_1 \to \mathrm{E}_2\) such that \(\gamma_2 = u \circ \gamma_1\) . In particular, we have \(x_2 = \gamma_2(e) = u(\gamma_1(e)) = u(x_1)\) . Moreover, for all \(g\) and \(h\) in \(\mathbf{G}\) , we have

\[
\varrho_ {2} (g) u \left(\gamma_ {1} (h)\right) = \varrho_ {2} (g) \gamma_ {2} (h) = \gamma_ {2} (g h) = u \left(\gamma_ {1} (g h)\right) = u \left(\varrho_ {1} (g) \gamma_ {1} (h)\right),
\]

and since the set of elements \(\gamma_{1}(h)\) for \(h \in G\) is total in \(E_{1}\) , this means that u is a G-morphism. By loc. cit., it is an isometric isomorphism if \((\varrho_{2}, x_{2})\) is also cyclic.

PROPOSITION 11. — Let \(\varphi \in \mathrm{Pos}_{1}(\mathrm{G})\) and \((\varrho, x)\) be a cyclic Hilbertian realization of \(\varphi\) . Then \(\varrho\) is an irreducible representation of G if and only if \(\varphi\) is an extreme point (EVT, II, p. 57, def. 1) of \(\mathrm{Pos}_{1}(\mathrm{G})\) .

Let E be the space of \(\varrho\) . Suppose first that \(\varrho\) is not irreducible. Let F be a closed subspace of E stable under \(\varrho\) , nonzero and different from E. Let us write \(x = x_1 + x_2\) , where \(x_1 \in \mathrm{F}\) and \(x_2 \in \mathrm{F}^\circ\) . We then have \(1 = \|x\|^2 = \|x_1\|^2 + \|x_2\|^2\) . The subrepresentation of E generated by \(x_1\) is contained in F, hence \(x_1 \neq x\) since \(x\) is a cyclic vector of \(\varrho\) , whence \(x_2 \neq 0\) . Similarly, one verifies that \(x_1 \neq 0\) .

For \(j = 1\) and \(j = 2\) , let \(\varphi_j\) be the continuous function on G such that

\[
\varphi_ {j} (g) = \frac {1}{\| x _ {j} \| ^ {2}} \langle x _ {j} | \varrho (g) x _ {j} \rangle
\]

for all \(g \in G\) . We have \(\varphi_{j} \in \mathrm{Pos}_{1}(G)\) (V, p. 443, example). Since \(\varphi = \|x_{1}\|^{2}\varphi_{1} + \|x_{2}\|^{2}\varphi_{2}\) , it suffices to check that \(\varphi_{1} \neq \varphi_{2}\) to show that \(\varphi\) is not an extreme point of \(\mathrm{Pos}_{1}(G)\) ; it suffices for that to prove that \(\varphi \neq \varphi_{1}\) .

Let us reason by contradiction and suppose that \(\varphi = \varphi_1\) . Since \(\langle x_1 | \varrho(g)x_2 \rangle = 0\) for all \(g \in G\) , we would have

\[
\frac {1}{\| x _ {1} \| ^ {2}} \langle x _ {1} | \varrho (g) x \rangle = \frac {1}{\| x _ {1} \| ^ {2}} \langle x _ {1} | \varrho (g) x _ {1} \rangle = \varphi_ {1} (g) = \varphi (g) = \langle x | \varrho (g) x \rangle
\]

for all \(g \in G\) , whence \(\langle x_{1} | y \rangle = \langle \|x_{1}\|^{2}x | y \rangle\) for every element y of the vector subspace of E generated by the elements \(\varrho(g)x\) for \(g \in G\) , therefore, for all \(y \in E\) , since x is a cyclic vector of \(\varrho\) . This would imply that \(x_{1} = \|x_{1}\|^{2}x\) is also a cyclic vector of \(\varrho\) , which is a contradiction, whence the assertion.

Suppose now that \(\varrho\) is irreducible and let us prove that \(\varphi\) is an extreme point of \(\mathrm{Pos}_1(\mathrm{G})\) . Let \(\varphi_1 \neq \varphi_2\) be elements of \(\mathrm{Pos}_1(\mathrm{G})\) and \(t_1, t_2 \in [0,1]\) such that \(t_1 + t_2 = 1\) and \(\varphi = t_1\varphi_1 + t_2\varphi_2\) . For \(j \in \{1,2\}\) , let \((\varrho_j, x_j)\) be a cyclic Hilbertian realization of \(\varphi_j\) , and let \(\mathrm{E}_j\) be the space of \(\varrho_j\) . Let \(x_3 = t_1^{1/2}x_1 + t_2^{1/2}x_2\) . Then \((\varrho_1 \oplus \varrho_2, x_3)\) is a Hilbert realization of \(\varphi\) . Since \((\varrho, x)\) is cyclic, there exists an isometric G-morphism \(u: \mathrm{E} \to \mathrm{E}_1 \oplus \mathrm{E}_2\) such that \(u(x) = x_3\) (prop. 10).

Let \(j = 1\) or \(j = 2\) . Since \(\varrho\) is irreducible, there exists \(\lambda_j \geqslant 0\) such that the G-morphism \(u_j = \mathrm{pr}_j \circ u\) from E into \(E_j\) satisfies

\[
\langle u _ {j} (y) \mid u _ {j} (y ^ {\prime}) \rangle = \lambda_ {j} \langle y \mid y ^ {\prime} \rangle
\]

for all \(y\) and \(y'\) in E (cor. 5 of V, p. 388, if \(u_j \neq 0\) , and one may take \(\lambda_j = 0\) otherwise). For every \(g \in G\) , it follows that

\[
\begin{array}{r} t _ {j} \varphi_ {j} (g) = \langle t _ {j} ^ {1 / 2} x _ {j} | \varrho_ {j} (g) (t _ {j} ^ {1 / 2} x _ {j}) \rangle = \langle u _ {j} (x) | \varrho_ {j} (g) (u _ {j} (x)) \rangle \\ = \langle u _ {j} (x) | u _ {j} (\varrho (g) x) \rangle = \lambda_ {j} \varphi (g). \end{array}
\]

Since \(\varphi_{j}(e) = \varphi(e) = 1\) , we deduce that \(\varphi_{j} = \varphi\) if \(t_{j} \neq 0\) . The hypothesis \(\varphi_{1} \neq \varphi_{2}\) therefore implies that \(t_{1}\) or \(t_{2}\) must be zero, which proves that \(\varphi\) is an extreme point of \(\mathrm{Pos}_{1}(\mathrm{G})\) .

Lemma 4. --- Let \(\varphi\in\mathrm{Pos}(\mathbf{G})\) . The function \(\varphi\) is bounded on \(\mathbf{G}\) and we have \(\|\varphi\|_{\infty}=\varphi(e)\) . Moreover, we have

\[
| \varphi (g ^ {- 1} h) - \varphi (h) | \leqslant \sqrt {2 \varphi (e) (\varphi (e) - \mathcal {R} (\varphi (g)))}\tag{1}
\]

for all \((g,h)\in \mathbf{G}\times \mathbf{G}\) .

Let \((\varrho, x)\) be a Hilbert realization of \(\varphi\) . We have \(\varphi(e) = \|x\|^2\) . For every \(g \in G\) , it follows therefore that \(|\varphi(g)| = |\langle x | \varrho(g)x\rangle| \leqslant \varphi(e)\) by the Cauchy-Schwarz inequality. This proves the first assertion.

For every \((g,h)\in \mathrm{G}\times \mathrm{G}\) , we have

\[
\varphi (g ^ {- 1} h) - \varphi (h) = \langle x | \varrho (g ^ {- 1} h) x \rangle - \langle x | \varrho (h) x \rangle = \langle \varrho (g) x - x | \varrho (h) x \rangle
\]

since \(\varrho\) is unitary, whence

\[
| \varphi (g ^ {- 1} h) - \varphi (h) | \leqslant \| \varrho (g) x - x \| \| \varrho (h) x \| \leqslant \varphi (e) ^ {1 / 2} \| \varrho (g) x - x \|.
\]

One concludes by observing that

\[
\| \varrho (g) x - x \| ^ {2} = 2 \| x \| ^ {2} - 2 \mathcal {R} (\langle x | \varrho (g) x \rangle) = 2 (\varphi (e) - \mathcal {R} (\varphi (g))).
\]

Remark. --- The sets \(\operatorname{Pos}(G)\) (resp. \(\operatorname{Pos}_{1}(G)\) and \(\operatorname{Pos}_{\leqslant1}(G)\) ) are convex self-adjoint subsets of the stellar algebra \(\mathcal{C}_{b}(G)\) . They are closed in the space \(\mathcal{C}_{b}(G)\) equipped with the topology of simple convergence. The set \(\operatorname{Pos}(G)\) is a convex cone with vertex 0 in the real Banach space \(\mathcal{C}_{b}(G)\) .

